\section{Limitations and Open Threats to Validity}

\begin{itemize}
  \item \textbf{Partial retail equivalence only.}  A shipped-executable Wine
  probe now agrees with the Linux reference for 2,000 frames and a 34-field
  raw projection on one replay.  It does not yet export the eleven canonical
  subsystem payloads, cover later stages or the replay matrix, prove that its
  input/timing interventions are noninterfering, or establish whole-state and
  whole-program equivalence.
  \item \textbf{The first microkernel is narrow.}  The integer PC24
  player-position transition matches 1,999 consecutive retail-derived steps,
  but it excludes dead/spawning writes and exceptional or subnormal arithmetic,
  and it does not model hitbox/orb/animation side effects.  Its dyadic rounding
  algorithm, reachable-domain checks, interior x87 control state, retail
  instruction binding, Rust/zkVM code binding, and projection noninterference
  are not yet machine-checked.
  \item \textbf{Floating point.}  The target uses x87 instructions including
  trigonometric and rounding operations, explicit DWORD/QWORD stores, and
  conversion helpers.  The static audit identifies wrapper word
  \code{0x027f}, while the retail probe observes ambient \code{0x007f} at all
  2,000 post-calc boundaries.  Neither result proves entry state or every
  transient per-instruction word, and the current PC53 SoftFloat campaign does
  not cover the observed PC24 ambient profile.  Hardware
  transcendental approximations also require a processor/emulator profile.
  Host \code{libm}, spill behavior, and intermediate precision may diverge by
  one or more bits and later alter collisions.  A pinned SoftFloat experiment
  agrees with hardware results and six exception bits for tested arithmetic,
  store, rounding, integer-conversion, and comparison cases, but it is neither
  formally verified nor complete x87 state semantics.  The comparison probe
  includes C0--C3 and selected NaNs, while a separate exact-byte
  \code{\_\_ftol2} probe covers full register results only for canonical finite
  signed-i64 inputs.  Load exceptions, arbitrary/noncanonical NaNs, remainder,
  exceptional helper inputs, site-specific reachable ranges/observations,
  address binding, and transcendental behavior remain open.  A focused ECL
  campaign tests exceptional/default classification but is not universal
  semantics.  A hash-bound ledger indexes 321 comparison/helper addresses; a
  derived manual
  source ledger proposes 68 omissions and nine retained calls, but every
  correspondence, slice disposition, and helper refinement is still open.  Its
  bounded scan also leaves EDX live at 42 stopping points.  The ECL quotient
  still lacks verified decoding, universal helper/classifier agreement,
  jump-target binding, and guest refinement; the eight point-score calls lack
  discharge of the dependent player-position artifact, collision-code binding,
  architectural masked-invalid/helper-path refinement, and caller binding,
  together with reachable difficulty and pre-score bounds.  The 95-instruction
  player audit and total clamp theorem reduce its invariant to non-NaN movement
  candidates plus initialization, writer-completeness, scheduling, binary/x87,
  correspondence, and guest obligations; they do not prove those premises.  A
  fixed-data ECL audit finds zero player-position outputs among 1,844 candidates
  and checks all 80 operations using the two handlers that bypass the readonly
  guard, but still lacks a
  proof of archive extraction, runtime decoding/dispatch, instruction
  immutability, non-ECL alias completeness, and artifact-to-Lean binding.
  The total score model covers finite, infinite, and NaN item coordinates
  conditionally; 2,081 Linux samples exercise only its finite branch and prove
  none of those bindings.  The complete address-indexed semantics, code
  binding, and any fixed-point refinement have not yet been implemented or
  proved.
  \item \textbf{Draw-chain effects.}  Headless mode skips drawing, but some draw
  callbacks write flags, ANM fields, laser positions, or other state that may be
  read later.  Each removal needs dependency analysis and corpus validation.
  \item \textbf{Incomplete corpus.}  The current samples cover Normal and
  Lunatic full runs but not all characters, shot types, Extra, retries, and
  adversarial stage-checkpoint mutations.
  \item \textbf{Reconstruction gaps.}  The matching project is highly complete
  but not identical in every function.  Unmatched code must be closed or
  explicitly excluded from the accepted statement.
  \item \textbf{Claim scope.}  A valid execution proof is not a proof of human
  play, authorship, or the absence of tool assistance.
  \item \textbf{Projection soundness.}  Equal selected-field digests imply
  equality only of that projection under the SHA-256 collision-resistance
  assumption.  Until coverage is frozen and justified, omitted state can still
  invalidate an observational-equivalence claim.  In particular, inactive
  Effect slots retain future-live values, and dynamic screen-effect calc jobs
  are not yet represented even though screen shake consumes the shared RNG.
  \item \textbf{Initializedness is semantic.}  Some reconstructed structures
  contain disabled members left uninitialized by the reference allocator.
  Hashing or constraining such bytes creates false distinctions and may import
  undefined host state into the statement.  Current ANM guards are tested but
  not yet justified by an exhaustive future-read proof.
  \item \textbf{Model-to-code binding.}  The planned Lean refinement is not yet
  connected to the authoritative C++ transition or a zkVM guest.  Differential
  agreement can expose mistakes in such a connection but cannot prove it.
\end{itemize}

The commercial game data is verified locally by digest and is not
redistributed.  Replay fixtures retain source links and redistribution status;
inputs whose sources prohibit reuse remain local.
